\chapter{Peritonitis}

\medskip
\noindent\textit{\textbf{Under review.} This chapter is an AI-assisted educational draft; it is not a substitute for professional medical advice.}

\section*{Definition and Overview}
inflammation or infection of the peritoneal lining, commonly from perforation, abdominal infection, surgery, or dialysis-related infection.

\section*{Clinical Features}
severe abdominal pain, guarding, fever, vomiting, bloating, reduced bowel movement, or systemic collapse.

\section*{When to Seek Medical Care}
Seek medical review for new, persistent, worsening, or function-limiting symptoms. Call 000 for severe breathing difficulty, collapse, sudden neurological change, severe chest or abdominal pain, or any rapidly worsening illness.

\section*{Diagnosis and Investigations}
urgent examination, blood tests and cultures, imaging, and analysis of peritoneal fluid when relevant.

\section*{Management}
urgent antibiotics, fluids, source control, and sometimes emergency surgery. Do not delay assessment for severe abdominal pain or collapse.

\section*{Complications}
Complications depend on the underlying cause and may include progression, effects on other organs, treatment adverse effects, or reduced quality of life. Early assessment can reduce preventable harm.

\section*{Prognosis and Self-Care}
Outcome depends on the cause, severity, complications, and response to treatment. Follow the care plan, keep appointments, avoid known triggers, and seek help if symptoms worsen.

\clearpage
